% Laget automatisk av elevark/lag_elevark.py fra opplegg/gravitasjonsloven.html. Endringer her blir overskrevet.
\documentclass[11pt]{article}
\usepackage{elevark}
\begin{document}
\elevarkhode{Fysikk 2 · Gravitasjon}{Newtons gravitasjonslov og gravitasjonsfelt – arbeidsark}{https://fysikklab.no/sim/satellittbaner.html}{Satellittbaner}

\section*{Oppgaver}

Oppgave 4–6 bruker Satellittbaner\footnote{\url{fysikklab.no/sim/satellittbaner.html}}. Bruk $γ = 6{,}67 \cdot 10^{-11}\,\mathrm{N\,m^{2}/kg^{2}}$, jordas masse $5{,}97 \cdot 10^{24}\,\mathrm{kg}$ og jordradien $6{,}37 \cdot 10^{6}\,\mathrm{m}$.

\begin{oppgaveliste}
\oppgave \fase{gjett}{Gjett}Jorda trekker på månen, og månen trekker på jorda. Hvilken kraft er størst?
\svar
\oppgave \fase{test}{Regn}Regn ut gravitasjonskraften mellom jorda og månen ($7{,}35 \cdot 10^{22}\,\mathrm{kg}$) når avstanden er $3{,}84 \cdot 10^{8}\,\mathrm{m}$.
\svar
\oppgave \fase{test}{Regn}Regn ut gravitasjonsfeltstyrken på bakken og i 400~km høyde. Hvor mange prosent av verdien på bakken er det?
\svar
\oppgave \fase{test}{Regn}Vis at $v = \sqrt{γM/r}$ i en sirkelbane, og regn ut farten og omløpstiden for romstasjonen i 400~km høyde. Sjekk med simuleringen.
\svar
\oppgave \fase{test}{Undersøk}Velg geostasjonær bane i simuleringen. Hvor lang er omløpstiden? Regn ut høyden selv fra $T = 23\,\mathrm{h}$ 56~min.
\svar
\oppgave \fase{forklar}{Forklar}Velg GPS-banen. Hvorfor har GPS-satellittene mindre fart enn romstasjonen, selv om de har lengre vei rundt?
\svar
\oppgave \fase{test}{Regn}Io går rundt Jupiter med baneradius $4{,}22 \cdot 10^{8}\,\mathrm{m}$ og omløpstid 1,769 døgn. Finn massen til Jupiter.
\svar
\oppgave \fase{utfordring}{Utfordring}Jorda går rundt sola med baneradius $1{,}496 \cdot 10^{11}\,\mathrm{m}$ på ett år. Finn massen til sola. Hvorfor kan vi ikke finne massen til jorda på denne måten?
\svar
\end{oppgaveliste}

\end{document}
